\documentclass[10pt,twocolumn]{article}

\usepackage[a4paper,margin=0.75in]{geometry}
\usepackage{times}
\usepackage{microtype}
\usepackage{amsmath}
\usepackage{amssymb}
\usepackage{booktabs}
\usepackage{listings}
\usepackage{xcolor}
\usepackage[hidelinks]{hyperref}

\lstdefinelanguage{LLVM}{
  morekeywords={define,declare,ret,br,switch,invoke,call,load,store,alloca,global,internal,private,external,trunc,zext,sext,and,or,select,phi,getelementptr,inbounds},
  sensitive=true,
  morecomment=[l]{;}
}
\lstset{
  basicstyle=\ttfamily\footnotesize,
  columns=fullflexible,
  breaklines=true,
  frame=single,
  language=LLVM
}

\title{Conservative Adaptive Bit-Width Narrowing for LLVM Memory Objects}
\author{Anonymous Artifact Draft}
\date{}

\begin{document}
\maketitle

\begin{abstract}
Integer values stored in memory often occupy fewer significant bits than their
declared source-level type. While scalar integer narrowing is a familiar
compiler optimization theme, narrowing memory objects is more delicate because
the transformation must preserve load/store semantics, pointer use structure,
and signedness restoration. This paper presents a conservative LLVM New Pass
Manager prototype, \textsc{AdaptiveWidthMemoryPass}, that narrows selected local
stack arrays and internal/private globals when all observed stores can be
classified by local bit-width inference. The pass rewrites memory object element
types, truncates stores, and restores narrowed loads with either zero-extension
or sign-extension according to a tracked extension family. Unsupported pointer
uses, mixed signed/unsigned value families, and unknown stored values are
rejected rather than speculated on. On the current fixture corpus, the prototype
identifies 28 memory-object candidates, transforms 20 of them, and records a
static summed per-element width reduction of 234 bits. On 32 C-derived
kernels, it transforms 26 memory-object candidates and exposes the expected
tradeoff between narrower memory and additional cast instructions. These results
are not presented as general benchmark claims; they establish a reproducible
artifact and a safety-first baseline for future memory bit-width optimization
experiments.
\end{abstract}

\section{Introduction}

Many programs store values in integer containers wider than the value range
required at a particular program point. Examples include masked counters,
small enumerations, byte-derived values stored in 16- or 32-bit arrays, and
sign-extended offsets later stored in wider integer memory. Narrowing such
memory objects can reduce static storage width and may enable downstream
optimizations. However, memory narrowing is more constrained than scalar
narrowing: every load and store must be rewritten consistently, pointer uses may
escape, and signedness-sensitive values must be restored with the correct
extension operation.

This work develops a small, inspectable LLVM pass that prioritizes correctness
and reproducibility over broad coverage. The prototype narrows only integer
arrays allocated on the stack and internal/private integer globals. It accepts
only direct pointer uses through getelementptr, load, and store. It performs
local stored-value inference over constants, casts, masks, simple select/phi/or
compositions, and LLVM's KnownBits information. If a candidate cannot be proven
safe and profitable under these rules, it is left unchanged and reported with a
reason.

The contributions of this draft are:
\begin{itemize}
  \item a conservative memory-object narrowing transformation for LLVM IR;
  \item an explicit extension-family safety policy for zero- and sign-restored
        loads;
  \item deterministic candidate classification and rejection logging; and
  \item a reproducible metrics harness that reports candidate coverage,
        transformation counts, static width deltas, and IR instruction overhead.
\end{itemize}

\section{Motivation}

Consider a local array whose stored values are masked before every store:

\begin{lstlisting}
%arr = alloca [2 x i16]
%masked = and i16 %x, 31
store i16 %masked, ptr %slot
%v = load i16, ptr %slot
\end{lstlisting}

The stored value needs only five bits. The pass can replace the array with
\texttt{[2 x i5]}, truncate the store, and zero-extend the load back to the
original load type:

\begin{lstlisting}
%arr.narrow = alloca [2 x i5]
%masked.narrow = trunc i16 %masked to i5
store i5 %masked.narrow, ptr %slot.narrow
%v.narrow = load i5, ptr %slot.narrow
%v.wide = zext i5 %v.narrow to i16
\end{lstlisting}

For signed byte-derived values, zero-extension is not equivalent. Therefore the
prototype tracks whether the stored-value family must be restored by
\texttt{zext} or \texttt{sext}. A candidate containing a mixture of both
families is rejected.

\section{Safety Model}

Let a candidate memory object have original element type $iN$ and proposed
narrow element type $iK$ where $K < N$. The transformation is allowed only if:

\begin{enumerate}
  \item the memory object is a supported stack array or internal/private global;
  \item all pointer users are recursively limited to supported GEP, load, and
        store operations;
  \item every store to the object has an inferred width at most $K$;
  \item every inferred stored value agrees on one extension family
        $E \in \{\mathrm{zero}, \mathrm{sign}\}$; and
  \item the candidate is profitable, i.e., $K < N$.
\end{enumerate}

Loads from the narrowed object are restored to $iN$ using the selected extension
family. Thus the intended local preservation condition is:

\[
  \mathrm{extend}_E(\mathrm{trunc}_K(v)) = v
\]

for each stored value $v$ that may be loaded from the object. The current
implementation checks this condition through syntactic and KnownBits-based
inference rules rather than whole-program analysis. If inference fails, the
candidate is treated as unknown and is not transformed.

\subsection{Invariant and Preservation Sketch}

For every transformed object $O$, the pass maintains the following local
invariant: every value stored through a supported pointer path to $O$ belongs to
one inferred extension family $E$ and has representable width at most $K$, the
new element width. Because all admitted loads and stores are rewritten through
the same rebuilt pointer graph, each transformed load observes a value previously
stored in truncated form. Since the inference phase requires
$\mathrm{extend}_E(\mathrm{trunc}_K(v)) = v$ for every admitted stored value,
restoring the narrowed load with $E$ preserves the original loaded value. Cases
outside this local argument--escaping pointers, unsupported casts, aggregate
loads, volatile/atomic behavior, unknown values, and mixed extension families--
are rejected before rewriting.

This is not a whole-program proof. It is a deliberately local preservation
argument for the subset of LLVM IR accepted by the implementation. The
deterministic rejection policy is therefore part of the safety model rather than
an engineering afterthought.

\section{Algorithm and Implementation}

The pass is implemented as an LLVM 18 New Pass Manager plugin. It exposes four
pipelines: full analysis/transformation, log-only analysis, stack-allocation-only
transformation, and global-only transformation.

\subsection{Candidate Discovery}

The pass scans functions for local stack integer arrays and scans the module for
internal/private integer globals. Scalar allocas, floating-point objects,
external globals, aggregate loads, and pointer escapes are rejected. A recursive
pointer-user walk admits GEP, load, and store users and records unsupported users
for diagnostic output.

\subsection{Width Inference}

The inference engine computes a required bit width and an extension family for
stored values. It currently recognizes:

\begin{itemize}
  \item integer constants;
  \item truncation, including C-style \texttt{trunc(and i32 \%x, mask)} forms;
  \item zero-extension and sign-extension chains;
  \item constant masks;
  \item bitwise OR;
  \item select and bounded phi nodes when all arms agree; and
  \item LLVM KnownBits as a conservative fallback.
\end{itemize}

For compositions such as select, phi, and OR, operands must agree on the same
extension family. This intentionally rejects mixed signed/unsigned candidates.

\subsection{IR Rewriting}

For a profitable candidate, the pass creates a replacement memory object with
the narrowed integer element type. It rebuilds GEPs using the new source element
type, rewrites stores by truncating stored values to the new type, rewrites loads
from the new type, and extends loaded values back to the original type. The old
load/GEP/object instructions are erased only after their replacements have been
installed and the resulting IR is verified by the test harness.

\subsection{Transformation Algorithm}

Algorithm~\ref{alg:narrowing} summarizes the pass at a high level. The pass is
structured as a sequence of filters: candidate discovery, use-graph validation,
stored-value inference, profitability testing, and finally IR rewriting. This
ordering is deliberate: it keeps unsound cases outside the rewriting phase and
makes each rejection observable in the diagnostic log.

\begin{figure}[t]
\begin{lstlisting}[language={}]
for object O in stack_arrays + internal_globals:
  if not integer_memory_object(O):
    reject(O, "non-integer")
    continue
  U = collect_supported_pointer_uses(O)
  if U has escaping or unsupported users:
    reject(O, "unsupported-user")
    continue
  W, E = infer_width_and_extension(stores(O))
  if inference failed:
    reject(O, "unknown-width")
    continue
  if W >= original_element_width(O):
    reject(O, "not-profitable")
    continue
  O2 = create_narrow_object(O, W)
  rewrite_geps_loads_stores(O, O2, E)
  erase_dead_old_ir(O)
\end{lstlisting}
\caption{High-level memory-narrowing algorithm.}
\label{alg:narrowing}
\end{figure}

\subsection{Pass Architecture}

The implementation separates four concerns: (1) candidate discovery, (2)
recursive pointer-use classification, (3) width and extension-family inference,
and (4) IR rewriting. The metrics harness consumes the same deterministic log
output used by the tests, so paper tables are generated from the artifact rather
than manually transcribed.

\section{Evaluation Methodology}

The artifact includes two evaluation layers. First, fixture-level tests exercise
individual transformation and rejection behaviors. Second, a benchmark-metrics
harness compiles small C kernels to LLVM IR with Clang 18, applies a light
\texttt{mem2reg} normalization step to remove irrelevant scalar spills, runs the
pass, verifies the transformed IR, and emits tab-separated and Markdown results.

Metrics include candidate counts, safety classifications, transformed globals
and allocas, rewritten loads and stores, static summed per-element bit savings,
LLVM IR instruction-count delta before and after LLVM default O2, and coarse
process-level pass wall time. The artifact also generates deterministic runtime
differential harnesses for each C kernel; baseline and transformed executables
are compiled and compared on checksum outputs. The wall-time and
instruction-count metrics are intended for artifact regression tracking, not as
rigorous performance claims.

\begin{table*}[t]
\centering
\caption{Fixture-level evaluation summary. The bit-savings column is a static per-element width delta summed across transformed candidates, not a measured runtime memory-footprint reduction.}
\label{tab:fixture-results}
\begin{tabular}{lrrrrrrrr}
\hline
Input & Cand. & Safe & Blocked & Unknown & Trans. G & Trans. A & Loads & Bits \\
\hline
candidates & 10 & 8 & 2 & 2 & 2 & 3 & 2 & 67 \\
transform\_alloca & 7 & 7 & 0 & 2 & 0 & 5 & 2 & 50 \\
inference & 6 & 6 & 0 & 0 & 0 & 6 & 5 & 59 \\
stack\_positive & 1 & 1 & 0 & 0 & 0 & 1 & 1 & 11 \\
globals.example & 2 & 2 & 0 & 0 & 2 & 0 & 2 & 36 \\
masked\_stack.example & 1 & 1 & 0 & 0 & 0 & 1 & 1 & 11 \\
rejected\_escape.example & 1 & 0 & 1 & 0 & 0 & 0 & 0 & 0 \\
\hline
Total & 28 & 25 & 3 & 4 & 4 & 0 & 13 & 234 \\
\hline
\end{tabular}
\end{table*}

\begin{table*}[t]
\centering
\caption{C-kernel benchmark metrics. Raw IR delta is measured immediately after rewriting; optimized IR delta is measured after LLVM default O2 on baseline and transformed IR.}
\label{tab:benchmark-results}
\begin{tabular}{lrrrrrrrr}
\hline
Benchmark & Cand. & Safe & Blocked & Trans. & R/W & Raw $\Delta$ & O2 $\Delta$ & Bits \\
\hline
alias\_pointer\_reject & 1 & 0 & 1 & 0 & 0 & +0 & +0 & 0 \\
alpha\_blend\_weights & 1 & 1 & 0 & 1 & 7 & +7 & +0 & 8 \\
checksum\_nibbles & 1 & 1 & 0 & 1 & 8 & +8 & +0 & 12 \\
embedded\_counters & 1 & 1 & 0 & 1 & 8 & +8 & +0 & 11 \\
fir4\_coefficients & 1 & 1 & 0 & 1 & 8 & +4 & +0 & 13 \\
global\_constants & 3 & 2 & 1 & 2 & 2 & +2 & +0 & 25 \\
global\_flag\_table & 3 & 3 & 0 & 3 & 3 & +3 & +0 & 42 \\
global\_lookup\_small & 1 & 1 & 0 & 0 & 0 & +0 & +0 & 0 \\
grayscale\_threshold & 1 & 1 & 0 & 1 & 8 & +8 & +0 & 8 \\
ipv4\_header\_fields & 1 & 1 & 0 & 1 & 8 & +8 & +0 & 8 \\
mixed\_signed\_reject & 1 & 1 & 0 & 0 & 0 & +0 & +0 & 0 \\
moving\_average4 & 1 & 1 & 0 & 1 & 8 & +8 & +0 & 6 \\
overflow\_uncertainty & 1 & 1 & 0 & 0 & 0 & +0 & +0 & 0 \\
packet\_flags & 1 & 1 & 0 & 1 & 8 & +8 & +0 & 11 \\
pixel\_channels & 1 & 1 & 0 & 1 & 8 & +8 & +0 & 8 \\
protocol\_version\_table & 1 & 1 & 0 & 1 & 8 & +8 & +0 & 11 \\
quantized\_filter & 1 & 1 & 0 & 1 & 6 & +6 & +0 & 9 \\
rejection\_cases & 3 & 2 & 1 & 0 & 0 & +0 & +0 & 0 \\
rgb565\_unpack & 1 & 1 & 0 & 1 & 6 & +6 & +0 & 10 \\
saturating\_sensor & 1 & 1 & 0 & 1 & 4 & +4 & +0 & 7 \\
select\_ranges & 1 & 1 & 0 & 1 & 4 & +4 & +4 & 10 \\
sensor\_bucket & 1 & 1 & 0 & 1 & 4 & +4 & +0 & 6 \\
signed\_delta\_pair & 1 & 1 & 0 & 1 & 4 & +4 & +0 & 8 \\
signed\_offsets & 1 & 1 & 0 & 1 & 4 & +4 & +0 & 8 \\
signed\_select\_offsets & 1 & 1 & 0 & 0 & 0 & +0 & +0 & 0 \\
signed\_state\_delta & 1 & 1 & 0 & 0 & 0 & +0 & +0 & 0 \\
small\_state\_machine & 1 & 1 & 0 & 1 & 6 & +6 & +0 & 13 \\
status\_flags & 1 & 1 & 0 & 1 & 6 & +6 & +0 & 13 \\
tcp\_flag\_decode & 1 & 1 & 0 & 1 & 12 & +12 & +0 & 15 \\
unsigned\_histogram & 2 & 2 & 0 & 1 & 5 & +5 & +0 & 9 \\
volatile\_reject & 1 & 0 & 1 & 0 & 0 & +0 & +0 & 0 \\
wide\_unknown\_reject & 1 & 1 & 0 & 0 & 0 & +0 & +0 & 0 \\
\hline
Total & 39 & 35 & 4 & 26 & 145 & +141 & +4 & 271 \\
\hline
\end{tabular}
\end{table*}


\section{Results}

Table~\ref{tab:fixture-results} shows the fixture-level result. Across seven
inputs, the pass identifies 28 candidates. It classifies 25 as safe under the
pointer-use policy and transforms 20 total memory objects: four globals and
sixteen stack arrays. The fixture corpus intentionally includes unknown-width
and escaped-pointer cases; these account for the three blocked candidates and
four unknown-width classifications. The resulting static summed per-element
width delta is 234 bits.

Table~\ref{tab:benchmark-results} shows the current C-kernel benchmark snapshot.
The suite includes unsigned mask extraction, pixel-channel packing, packet flag
decoding, signed offsets, global constants, select-based ranges, and deliberate
rejection cases. Across 32 kernels, the pass sees 39 memory-object
candidates, classifies 36 as safe under the use policy, and transforms 26
objects. The transformed C kernels increase static IR instruction count by 141
assignments because the pass inserts explicit truncation and extension
instructions; after LLVM default O2 this overhead falls to four assignments.
Runtime differential tests compile baseline and transformed versions of all
32 benchmark harnesses and compare deterministic checksum outputs. All
32 currently match, providing a practical behavioral sanity check in
addition to LLVM verifier checks. This does not replace formal proof or large
input-space testing, but it strengthens the artifact-level evidence for the
current benchmark suite.

\section{Limitations}

The current prototype is intentionally narrow. It does not handle arbitrary
pointer aliasing, external globals, scalar allocas, bitcast/address-space-cast
rewrite cases, aggregate loads, volatile or atomic semantics, interprocedural
effects, or broad arithmetic range inference. The current evaluation uses small
fixtures and kernels rather than external benchmark suites. Therefore the
reported numbers establish correctness and reproducibility of the artifact, not
general application-level memory or performance improvements.

\section{Related Work}

The implementation builds on LLVM's modular pass infrastructure
\cite{lattner2004llvm} and uses LLVM IR semantics as the transformation target.
Its inference style is related to classic data-flow and abstract-interpretation
ideas \cite{cousot1977abstract}, and to compiler propagation analyses such as
conditional constant propagation \cite{wegman1991sccp}. The use of LLVM
KnownBits connects the prototype to existing value-tracking infrastructure.
Unlike scalar-only simplifications, this work focuses specifically on narrowing
memory object element types while preserving load/store behavior.

\section{Future Work}

Future work should scale the metrics harness to larger benchmark suites,
measure optimized machine-code effects after downstream LLVM passes, support
additional safe pointer-use patterns, extend arithmetic range inference, and add
runtime differential testing for transformed kernels. Another important
direction is comparing memory-object narrowing against scalar bit-width
optimizations to determine when memory rewriting provides unique benefits.

\section{Conclusion}

This draft presented a conservative LLVM prototype for adaptive bit-width
narrowing of selected memory objects. The pass combines local stored-value
inference, explicit signedness restoration, conservative pointer-use filtering,
and deterministic reporting. The current results show that the implementation
can safely transform targeted fixture and C-derived kernels while preserving a
clear rejection policy. The artifact is now positioned for broader benchmark
evaluation and iterative paper development.

\bibliographystyle{plain}
\bibliography{references}

\end{document}
